% Chapter 4. Method and Conventions.
% Every number in this chapter comes from verify/registry.json (keys named in the
% text), from verify/recompute_25m0328.py, or from the committed model output.

\chapter{Method and Conventions}
\label{ch:conv}

Every experiment in Chapters~5 and~6 compares two converged model states and reports
how they differ. A difference can be measured in more than one way, and two honest
calculations from the same output can give different numbers if they make different
choices. This chapter states every choice that was made, so that a reader who repeats
a calculation gets the number printed here: how numbers are written
(Section~\ref{sec:conv-numbers}), the six conventions with a worked example each
(Section~\ref{sec:conv-six}), the fixed ranges used throughout
(Section~\ref{sec:conv-windows}), how the size of a change in the whole radiation field
is measured (Section~\ref{sec:conv-field}), the response band
(Section~\ref{sec:conv-band}) and the reproducibility floor
(Section~\ref{sec:conv-floor}).

The same six conventions are also written out in the code that computes every number,
\texttt{verify/recompute\_25m0328.py}: the first five at the top of the file, and the
sixth as the value of the key \texttt{res.e10.SO.inband}. Where a recomputation disagrees with a number in
this report, the first thing to check is whether it follows these conventions.

% ===========================================================================
\section{How numbers are written in this report}
\label{sec:conv-numbers}

The registry, the file \texttt{verify/registry.json}, holds every quantity computed from
the model, each under its own name, called a \emph{registry key}; an example is
\texttt{e6.A1.meanz}. Chapter~8 explains how to read these names. Where this report prints such a number, the key is
named next to it or at the start of the paragraph.

A printed number carries only as many digits as the text needs. It is correct when the
registry value, rounded to the same number of significant digits, gives the printed
number. For example, \texttt{e6.A1.meanz} holds the response centroid to sixteen digits,
and the text prints it as 82.044~km. The project's checker,
\texttt{verify/check\_documents\_25m0328.py}, applies this rule against the whole
registry to every number written with three or more decimal places, in scientific
notation, or as a percentage with two or more decimal places.
Percentages may be stored either as percentages or as fractions; 1.4484\% and a stored
0.014484 are the same number.

The checker can show that a printed number exists in the registry. It cannot show that it
was taken from the right key. That second check is done by hand, by recomputing the
sentence from the key it names, and Chapter~8 gives an example where it mattered.

% ===========================================================================
\section{The six conventions}
\label{sec:conv-six}

Table~\ref{tab:conventions} lists the six conventions. Each is explained below with a
worked example.

\begin{table}[htbp]
  \centering
  \caption[The six conventions used for every number in this report]{The six conventions used for every number in this report.}
  \label{tab:conventions}
  \footnotesize
  \setstretch{1.0}
  \begin{tabularx}{\textwidth}{>{\raggedright\arraybackslash}p{0.7cm}
                                >{\raggedright\arraybackslash}p{4.2cm}
                                >{\raggedright\arraybackslash}X}
    \toprule
    No. & Convention & What it means \\
    \midrule
    1 & Relative change & The difference is divided by the larger of the two values,
      not by the control. \\
    2 & Means over a field & Every element counts equally. A \emph{bounded} mean
      leaves out the elements where both values are zero. \\
    3 & Integrals over wavelength & VULCAN's own rule: two bin widths, and half weight
      at the two ends of every range. \\
    4 & Integrals over altitude & Layer thickness in centimetres, never kilometres. \\
    5 & Two photolysis measures & $\Delta J$ (the rate coefficient) and $\Delta(J n)$
      (the rate) are different quantities and are never mixed. \\
    6 & Reproducibility floor & In-band changes in abundance smaller than 2.5278\% are
      not resolved. \\
    \bottomrule
  \end{tabularx}
\end{table}

% ===========================================================================
\subsection{Relative change divides by the larger value}
\label{sec:conv-rel}

When a quantity $X$ has the value $X_{\mathrm{r}}$ in the reference run and
$X_{\mathrm{c}}$ in the case being tested, its relative change is
\begin{equation}
  R_X = \frac{|X_{\mathrm{c}} - X_{\mathrm{r}}|}
             {\max\big(|X_{\mathrm{r}}|,\,|X_{\mathrm{c}}|\big)} .
  \label{eq:rel}
\end{equation}
The denominator is the larger of the two magnitudes. This keeps $R_X$ between 0 and 1,
or between 0 and 100 per cent, whatever the two values are, and it is defined even when
the reference value is zero. For a small change the larger value and the reference
value are almost equal, and the two possible definitions agree. They differ only when
the change is large, which is exactly where a percentage needs care.

\textbf{Worked example.} In Experiment~2 the mixing ratio of OSSO (S\sub{2}O\sub{2}) at
33~km rises from $4.012\times10^{-42}$ with the absorber off to $3.0845\times10^{-41}$
with it on (registry keys \texttt{e2.S2O2.ref\_at\_level} and
\texttt{e2.S2O2.case\_at\_level}). The change, $2.683\times10^{-41}$, divided by the
larger value gives $R = 0.8699$ (\texttt{e2.S2O2.maxrel}). Divided by the control
instead it would give about 669 per cent, a number that says more about how tiny the
control value is than about the change. A second case shows why the rule is needed at
all: ClS\sub{2} at 97~km is exactly zero with the absorber off and
$6.9437\times10^{-27}$ with it on (\texttt{e2.ClS2.case\_at\_level}). A change relative
to the control is then undefined, while Equation~\eqref{eq:rel} gives exactly 1.

Both values in this example are far below anything that could be observed. That is why
Chapters~5 and~6 always give an absolute value next to a relative change for trace
species.

% ===========================================================================
\subsection{Means over a field}
\label{sec:conv-mean}

The model saves its fields as arrays: the actinic flux, for example, has one value for
each of 57 layers and 1670 wavelength bins, 95190 values in all. Where a mean is
reported over such an array, every element counts equally; no element is weighted by
layer thickness or bin width. The one exception is the \emph{bounded} mean of
Equation~\eqref{eq:rel}, which is taken only over elements where the denominator is not
zero:
\begin{equation}
  \overline{R} = \frac{1}{M} \sum_{\text{elements with } \max(|X_{\mathrm{r}}|,|X_{\mathrm{c}}|) > 0} R_X ,
  \label{eq:boundedmean}
\end{equation}
where $M$ is the number of those elements. An element where both runs give exactly zero
has no change to report, and counting it as a zero would dilute the mean by an amount
that depends only on how much of the grid is dark.

\textbf{Worked example.} Between the runs with the absorber off and on (Experiments~1
and~2), the actinic flux is non-zero in only 37185 of its 95190 elements; the rest are
short wavelengths that are absorbed completely before they reach that layer, and the
lower the layer, the more of the spectrum is dark (\texttt{e3fp.bnd.A1.aflux.n}). The bounded mean relative change over the 37185 lit
elements is 3.7078\% (\texttt{e3fp.bnd.A1.aflux.mean\_pct}). The mean over all 95190
elements, with each dark element counted as zero, is 1.4484\%
(\texttt{e2.aflux.meanrel}). The two differ by exactly the share of lit elements, and
both are correct for what they say. This report states which one it uses.

% ===========================================================================
\subsection{Integrals over wavelength}
\label{sec:conv-quad}

The photolysis rate coefficient of Equation~\eqref{eq:J} is an integral over
wavelength, which the model evaluates as a sum over its bins. The bins are 0.1~nm wide
from 96 to just below 240~nm (the first 1440 bins) and 2~nm wide from 240 to 698~nm
(\texttt{nom.base.v.dbin1} and \texttt{.dbin2}). VULCAN adds up each of the two parts
with the trapezoidal rule, which gives the two end bins of each part half weight. For a
range of bins from $a$ to $b$ with width $\Delta\lambda$,
\begin{equation}
  \int_{a}^{b} f(\lambda)\,\mathrm{d}\lambda \approx
  \sum_{i=a}^{b} f_i\,\Delta\lambda \;-\; \tfrac{1}{2}\big(f_a + f_b\big)\,\Delta\lambda ,
  \label{eq:quad}
\end{equation}
where $f_i$ is the integrand in bin~$i$. Every band integral in this report, for
example the share of a response that falls between 320 and 340~nm, uses the same rule
at its own two edges.

\textbf{Worked example.} A band from 320 to 340~nm holds 11 bins of 2~nm, centred at
320, 322, and so on up to 340~nm. For an integrand equal to 1 the true integral is 20~nm.
A plain sum of the 11 bins gives 22~nm, a tenth too much; Equation~\eqref{eq:quad}
removes half a bin at each end and gives 20~nm. With this rule, recomputing the stored
photolysis rate coefficients from the stored actinic flux and cross sections gives back
the model's own values: the difference is exactly zero in 102 of the 110 checks made
(55 photolysis branches in two runs), and very small in the rest (registry keys ending
\texttt{recon\_maxerr}).

% ===========================================================================
\subsection{Integrals over altitude}
\label{sec:conv-col}

Many model quantities are per cubic centimetre: number densities, and reaction rates in
molecules per cubic centimetre per second. Integrating such a quantity over altitude
gives a column, per square centimetre, only if each layer thickness is in centimetres:
\begin{equation}
  C = \sum_{j} r_j\,\Delta z_j , \qquad \Delta z_j \text{ in cm} .
  \label{eq:col}
\end{equation}
A thickness of 2~km is $2\times10^{5}$~cm, so taking it in kilometres makes the result
exactly $10^{5}$ times too small and leaves its unit meaningless.

\textbf{Worked example.} In Experiment~7, moving the absorber upward changes the contribution
that the photolysis of SO\sub{2} (reaction 781 of the network) makes to the SO\sub{2}
budget. Integrated over the response band with the thickness in centimetres, the change
is $+8.2352\times10^{8}$ molecules per square centimetre per second
(\texttt{e7.7a.SO2.R781.col}). Taking the thickness in kilometres instead gives
8235.1671, the same digits $10^{5}$ times smaller (\texttt{e7.7a.SO2.R781.col\_km}).
The registry keeps both so that the size of the error can be seen; only the first is
used in this report.

% ===========================================================================
\subsection{Two photolysis measures that must not be mixed}
\label{sec:conv-J}

A change in photolysis can be reported in two ways. The change in the rate
coefficient, $\Delta J$, is the change in the probability per second that one molecule is
broken up. It depends only on the sunlight reaching that height, so it is a purely
radiative measure. The change in the rate, $\Delta(J n)$, where $n$ is the number
density of the molecule, is the change in how many molecules are broken up per cubic
centimetre per second. It also depends on how much of the molecule is present. Because
most molecules become more abundant lower down while sunlight becomes weaker, the
second measure usually peaks lower in the atmosphere than the first.

Experiment~6 reports $\Delta(J n)$. Experiments~7 to~10 report $\Delta J$ when they
describe photolysis; the change in a reaction's contribution to a budget, as in
Experiments~5 and~7, is a rate and so of the $\Delta(J n)$ kind.

\textbf{Worked example.} For the photolysis of SO\sub{3} in Experiment~9 (Chapter~6),
the largest change of the rate coefficient inside the response band is 0.05898\%, at
81~km (\texttt{dJ.e9.SO3} and \texttt{dJz.e9.SO3}). The largest change of the rate is
0.4603\%, at 73~km (\texttt{dJy.e9.SO3} and \texttt{dJyz.e9.SO3}). The two measures
describe the same run and differ in size by a factor of about eight and in height by
8~km. A number taken from one and placed in a sentence about the other would look like
a disagreement in altitude where there is none.

% ===========================================================================
\section{Fixed ranges used throughout}
\label{sec:conv-windows}

Four ranges are used again and again, and each has the same limits everywhere in
\texttt{verify/recompute\_25m0328.py}, so that no analysis picks its own. The
\emph{absorber grid}, 300 to 698~nm, is where the absorber term acts: the absorber table
starts at 300~nm and the model grid ends at 698~nm. The \emph{enhancement band}, 320 to
400~nm, is the near-ultraviolet feature the absorber has to supply
(Chapter~\ref{ch:intro}). The \emph{screening window}, 300 to 620~nm, is the range
searched for each candidate's strongest absorption in the screens of Experiments~8
and~12. The \emph{response band}, 69 to 93~km, is where the model's radiation responds to
the absorber (Section~\ref{sec:conv-band}). Experiment~4 also uses narrower analysis
windows, among them 420 to 620~nm; these are deliberate divisions of the spectrum for
that experiment, not errors in the ranges above.

% ===========================================================================
\section{The size of a change in the radiation field}
\label{sec:conv-field}

Most radiative results in Chapters~5 and~6 compare a whole actinic flux field, 57 layers
by 1670 wavelength bins, with a reference field, and report the comparison as one
number: the size of the change as a percentage of the size of the reference field,
\begin{equation}
  \Delta F = 100\,
  \frac{\Big[\sum_{j}\sum_{i} w_i\,\big(F_{\mathrm{c}}(j,\lambda_i)
        - F_{\mathrm{r}}(j,\lambda_i)\big)^{2}\Big]^{1/2}}
       {\Big[\sum_{j}\sum_{i} w_i\,F_{\mathrm{r}}(j,\lambda_i)^{2}\Big]^{1/2}}
  \ \text{per cent},
  \label{eq:field}
\end{equation}
where $F_{\mathrm{c}}$ and $F_{\mathrm{r}}$ are the actinic flux in the case and in the
reference, $j$ runs over every layer, $i$ over every wavelength bin, and $w_i$ is the
width of bin~$i$: 0.1~nm below 240~nm and 2~nm above (the array \texttt{quad} in
\texttt{verify/recompute\_25m0328.py}). Each layer counts once, whatever its thickness.
Squaring the differences lets the largest changes dominate, and dividing by the size of
the reference field turns the result into a share of the light that is there. The
function \texttt{radiative()} in the same file computes it, and the registry holds it
under keys ending in \texttt{L2pct}.

\textbf{Worked example.} Moving the absorber 18~km upward (Experiment~7A), with the
absorber in its usual place as the reference, gives $\Delta F = 2.2498\%$
(\texttt{e7a.L2pct}); moving it 18~km down gives 0.3067\% (\texttt{e7b.L2pct}). The
spectral substitutions of Experiments~8 and~11 use the same formula with the baseline of
Experiment~8 as the reference.

One quantity uses the same sums with a different denominator. The radiation floor of
Section~\ref{sec:conv-floor} divides the difference between two runs by the change that
one of them makes, not by the reference field, so it is a share of a change, not a share
of the light. This is stated wherever it is used.

% ===========================================================================
\section{The response band}
\label{sec:conv-band}

The response band answers one question: at what altitudes does the model's radiation
actually respond to the absorber? Experiment~6 answered it by comparing the runs with
the absorber on and off. At each layer $j$ it adds up the change in actinic flux over
all wavelengths, weighting each bin by its width:
\begin{equation}
  w_j = \sum_{i} \big|F_{1}(j,\lambda_i) - F_{0}(j,\lambda_i)\big|\,\Delta\lambda_i ,
  \label{eq:weight}
\end{equation}
where $F_1$ and $F_0$ are the actinic flux with the absorber on and off. The quantity
$w_j$ is the size of the radiative response at layer $j$, and its distribution over
altitude shows where the response lives. It is purely radiative, so the convention of
Section~\ref{sec:conv-J} does not arise here.

The weighted mean altitude of this distribution, the response centroid, is 82.044~km
(\texttt{e6.A1.meanz}). A quarter of the response lies below 77~km, half below 81~km and
three quarters below 87~km (\texttt{e6.A1.q25}, \texttt{.q50} and \texttt{.q75}). The
\emph{response band} is the narrowest continuous range of altitudes that holds at least
90 per cent of the response. It runs from 69 to 93~km and holds 91.370\% of it
(\texttt{e6.interval.90.lo}, \texttt{.hi} and \texttt{.contained}). Table~\ref{tab:band}
shows how the range would change if a different share were chosen. The lower edge stays
at 69 or 71~km for every share from 68 to 95 per cent, so the band does not depend
sensitively on the choice of 90.

\begin{table}[htbp]
  \centering
  \caption[The narrowest altitude range that holds a given share of the radiative response (Experiment~6)]{The narrowest altitude range that holds a given share of the radiative
    response (Experiment~6). Registry keys \texttt{e6.interval.<share>.lo}, \texttt{.hi}
    and \texttt{.contained}.}
  \label{tab:band}
  \footnotesize
  \setstretch{1.0}
  \begin{tabular}{cccc}
    \toprule
    Share asked for & Range (km) & Width (km) & Share actually held \\
    \midrule
    50\,per cent & 73 to 83 & 10 & 54.567\% \\
    68\,per cent & 71 to 85 & 14 & 68.163\% \\
    80\,per cent & 71 to 89 & 18 & 81.479\% \\
    90\,per cent & 69 to 93 & 24 & 91.370\% \\
    95\,per cent & 69 to 97 & 28 & 95.240\% \\
    99\,per cent & 65 to 103 & 38 & 99.138\% \\
    \bottomrule
  \end{tabular}
\end{table}

The band is a property of the atmosphere, not of the absorber's spectrum. Experiment~4
repeated the run six times with the absorber confined to six different wavelength
ranges. The response centroid of those six runs lies between 81.989 and 82.100~km
(\texttt{e6.B5.meanz} and \texttt{e6.B6.meanz}), within a fraction of one model layer of
the full run.

The band lies above most of the absorber's own particles: 90.194\% of the mode~1
profile is between 58 and 69~km (\texttt{pnom.cloud}). The response to the absorber is
therefore felt mainly higher up than the absorber sits. Chapter~5 shows the altitude
profile of the response in full (Experiment~6). Every candidate test from Experiment~7
onward asks how much of a candidate's effect falls inside this band.

% ===========================================================================
\section{The reproducibility floor}
\label{sec:conv-floor}

A model run stops when the steady-state tests of Section~\ref{sec:how} pass, and those
tests allow mixing ratios to keep drifting slightly. Two runs that ought to give the same
answer can therefore end a little apart. The size of that disagreement sets the
smallest change in abundance that can be called a result.

Experiments~9 and~10 measured it, although they were not designed to. The two runs put
two different absorber spectra on the same low vertical profile, and Chapter~6 shows why
that made them radiatively identical: the difference between their actinic flux fields
is only 0.0017\% of the change that Experiment~9 makes (\texttt{e10.fielddiff\_pct}),
and that change is itself only 0.3245\% of the field (\texttt{e9.L2pct},
Equation~\eqref{eq:field}). Any difference between their chemistry must then be
numerical, not physical. Yet, measured against the same baseline, the largest change in
the SO mixing ratio inside the response band is 0.7384\% in one run and 2.5278\% in the
other (\texttt{res.e9.SO.inband} and \texttt{res.e10.SO.inband}). For SO\sub{3} the two
are 0.3488\% and 1.4288\% (\texttt{res.e9.SO3.inband} and \texttt{res.e10.SO3.inband}).

\textbf{The rule.} A change in the abundance of a species inside the response band that
is smaller than 2.5278\% is not a measurement. It can be reported only as an upper bound:
the true change is no larger than this, and may be zero. The convergence tolerance of
0.01 in \texttt{vulcan\_cfg.py} might suggest a floor of about 1 per cent; the measured floor is two and a half times larger. It comes from a single pair of runs, so
it is the smallest value the floor can have, not a full description of the scatter.

\textbf{What it bounds, and what it does not.} The floor applies to abundances, and so
also to any quantity that contains an abundance, such as the photolysis rate
$\Delta(J n)$. It does not apply to purely radiative quantities such as the actinic flux
or $\Delta J$: the same pair of runs agrees on those to 0.0017\% of the change itself,
so a radiative difference well below one per cent is still resolved. Of the 45 in-band changes in
the sulfur reservoirs held for Experiments~7 to~11 (keys beginning \texttt{res.}, which
divide by the reference run rather than as in Convention~1), only three reach the floor:
SO and SO\sub{3} in Experiment~7A, 8.404\% and 4.718\%, and the SO value that defines
the floor itself. The rest are upper bounds.

\textbf{Every run converged.} The floor is measured between converged runs, so it
matters that every run met the steady-state tests of Section~\ref{sec:how}.
Table~\ref{tab:conv} gives the two tested quantities for the main runs of Chapters~5
and~6; all pass, as do the control and the runs of Experiments~3 and~4B.

\begin{table}[htbp]
  \centering
  \caption[The steady-state tests at the end of each main run]{The steady-state tests at the end of each main run, with limits 0.01 and
    $10^{-5}$~s$^{-1}$. Keys \texttt{conv.<run>.longdy} and \texttt{conv.<run>.longdydt}.}
  \label{tab:conv}
  \scriptsize
  \setstretch{1.0}
  \setlength{\tabcolsep}{3pt}
  \begin{tabular}{lccc@{\hspace{1.2em}}lccc}
    \toprule
    Run & Key & \texttt{longdy} & \texttt{longdydt} & Run & Key & \texttt{longdy} & \texttt{longdydt} \\
    \midrule
    Exp.~2, on & \texttt{a1} & 0.0099963 & $6.991\times10^{-11}$ & Exp.~8C, S\sub{3} & \texttt{e8c} & 0.0099686 & $7.006\times10^{-11}$ \\
    Exp.~7A, up & \texttt{e7a} & 0.0098794 & $6.949\times10^{-11}$ & Exp.~9 & \texttt{e9} & 0.0098368 & $6.916\times10^{-11}$ \\
    Exp.~7B, down & \texttt{e7b} & 0.0099450 & $6.992\times10^{-11}$ & Exp.~10 & \texttt{e10} & 0.0097778 & $6.841\times10^{-11}$ \\
    Baseline, 8 to 11 & \texttt{b8} & 0.0099437 & $6.993\times10^{-11}$ & Exp.~11A, SSCl\sub{2} & \texttt{e11a} & 0.0099988 & $7.029\times10^{-11}$ \\
    Exp.~8A, ClS\sub{2} & \texttt{e8a} & 0.0099390 & $6.987\times10^{-11}$ & Exp.~11B, ClSSCl & \texttt{e11b} & 0.0099701 & $7.012\times10^{-11}$ \\
    Exp.~8B, OSSO & \texttt{e8b} & 0.0098173 & $6.901\times10^{-11}$ & & & & \\
    \bottomrule
  \end{tabular}
\end{table}


